\chapter*{Abstract}
\addcontentsline{toc}{chapter}{Abstract}

This diploma project is devoted to the development of a smart plant condition monitoring system that combines Internet of Things technologies, web application components, database storage, and an explainable artificial intelligence module. The practical motivation of the work is connected with the fact that indoor plant care is often based only on visual observation. A user may notice a problem only when the plant already shows visible stress, while environmental indicators can change much earlier. For this reason, automatic monitoring of soil moisture, air temperature, air humidity, and light intensity can support more timely and understandable plant care decisions.

The aim of the project is to design and implement a functional system that collects plant-related environmental data, stores it in a database, analyzes the current condition of the plant, generates alerts and recommendations, and displays the result through a web dashboard. The implemented system represents a functional diploma prototype with end-to-end functionality. It includes an Arduino Uno-based sensor node, a Python serial gateway, a FastAPI backend, a PostgreSQL database, and a React web dashboard. The Arduino Uno-based sensor node collects environmental measurements and sends them as JSON-formatted USB serial messages. The Python serial gateway reads these messages, checks and normalizes the values, and forwards structured telemetry to the backend.

The backend is responsible for the main data processing and storage logic. It validates incoming readings, stores measurements in PostgreSQL, manages users, plants, sensors, alerts, recommendations, notifications, and system logs, and provides processed information to the frontend. The frontend presents plant profiles, current sensor readings, historical values, condition interpretation, alerts, and recommendations. User access is protected through JSON Web Tokens, while sensor data ingestion is protected through a device token mechanism for registered devices. This separation allows the system to distinguish between ordinary user access and device-level telemetry submission.

The analytical part of the system is implemented as a hybrid module. The rule-based component is used as the primary explainable layer. It evaluates current sensor values and recent trends against predefined ranges and calculates condition status, health score, confidence estimate, risk factors, and a human-readable explanation. The machine learning component is used as an additional support layer. A supervised machine learning classifier based on \texttt{RandomForestClassifier} provides an additional condition prediction using moisture, temperature, humidity, light, and short-term trend features as input data. The model classifies plant condition into three classes: normal, attention, and critical. Since it is trained on synthetic scenario data generated for the project, its result is treated as decision support and does not replace the primary rule-based logic.

The scope of the system is limited to environmental plant condition monitoring and care recommendation support. The developed prototype does not perform biological diagnosis, plant disease detection, computer vision analysis, deep learning, actuator control, or industrial greenhouse automation. This limitation is important because the system uses environmental sensor readings and synthetic scenario data rather than long-term real agronomic datasets.

The developed prototype was verified through database migration checks, automated backend tests, frontend build validation, API checks, protected sensor ingestion checks, dashboard response validation, alert retrieval, and simulator-based demo flows. The results show that the system can receive telemetry, store measurements, analyze plant condition, generate recommendations and alerts, and visualize the processed data through the dashboard. The project demonstrates the practical feasibility of combining IoT integration and AI-supported analysis in a web-based plant monitoring solution within the scope of a diploma project.

\noindent\textbf{Keywords:} smart plant monitoring, Internet of Things, Arduino Uno, Python serial gateway, FastAPI, PostgreSQL, React, rule-based analysis, Random Forest, dashboard, alerts